% Replace the statement and proof of night:P1.18. The second splitting's
% numerical translation belongs in the check record, not the paper.
% Keep every old label/tag below; labels on computational displays remain
% attached to the corresponding mathematical input or conclusion.
\begin{theorem}[The minimum of the height function]
\hypertarget{tag-02QR}{}\label{night:P1.18}
\StatementRecord{night:P1.18}{proved}{02QR}{\SeriesRef{night:P1.15}, \SeriesRef{night:P1.8}}
The function $h_b$ has its unique minimum at $\lambda=\infty$.
For the normalized fundamental solution $H(1)=I$ of the algebraic
connection and its regular infinity limit $H_\infty$, this minimum is
\[
 h_{b,\min}=-\frac{96}{625\pi}\operatorname{Re}(XH_\infty)_{03}.
\]
Its negative $c_*=-h_{b,\min}$ satisfies the exact bounds
\[
 1.094540250311402780116268263054945703725<c_*
 <1.094540250311402780116268263054945703726.
 \hypertarget{tag-03KZ}{}\label{layout:V:fresh-bounds}
\]
The full parameter domain is $\mathcal U=\{\xi:\operatorname{Im}\xi>c_*\}$.
\end{theorem}
\begin{proof}
Put $a=256/3125$ and $c_0=96/(625\pi)$. We use the rational
connection of \cref{night:P1.15}, its integral matrices $A,B$,
and the normalization $Y=XH$, $H(1)=I$, $Xe_0=e_0$.
The finite-root gauges are
\[
 \lambda=u^4,\quad D_0=\operatorname{diag}(1,u^6,u^3,1,u^{-3},u^{-6}),
 \qquad
 \lambda=u^{-5},\quad D_\infty=
 \operatorname{diag}(1,u^{-8},u^{-4},1,u^4,u^8).
 \hypertarget{tag-03KX}{}\label{layout:V:fresh-gauges}
\]
For either gauge the regular connection is
$D^{-1}\lambda'\Gamma_6D+D^{-1}D'$.
Let $U_0,U_\infty$ be its solutions with value $I$ at zero.
The deck characters $T_0,T_\infty$ give the exact based
monodromies $U_0(1)^{-1}T_0^{-1}U_0(1)$ and
$U_\infty(1)^{-1}T_\infty^{-1}U_\infty(1)$.
The unique common intertwiner fixing $e_0$ is $X$, by
\cref{night:P1.15}. Hence the regular limiting frames are
$F_0=XU_0(1)^{-1}$ and $F_\infty=XU_\infty(1)^{-1}$.

The height is obtained from a real linear system. Write
$q=(0,v_2,v_1,-i,0,0)^t$ and impose that the lower five
coordinates of $Yq$ be real. The four real unknowns are the
real and imaginary parts of $v_2,v_1$; existence and uniqueness
are those of \cref{night:P1.8}. In these coordinates
$h_b=c_0\operatorname{Im}(Yq)_0$.
Every monodromy is real and fixes $e_0$, so the imaginary top
coordinate is unchanged when the lower coordinates are real.
This proves that the calculation gives one function on the base.
The real normalized conjugation matrix $K=X\bar X^{-1}$ satisfies
$Y(\bar\lambda)=K\overline{Y(\lambda)}$ and $Ke_0=e_0$;
using $-\bar q$ gives $h_b(\bar\lambda)=h_b(\lambda)$.

We give the enclosure method and the resulting whole-domain
inequalities. For an analytic right connection $G$ on a disc
$|z-z_0|\le R$, with a submultiplicative norm bound $\|G\|\le M$,
its normalized solution satisfies $\|U\|\le e^{MR}$.
Gronwall's estimate and Cauchy's coefficient formula therefore
bound the tail after degree $N$, on $|z-z_0|\le r<R$, by
\[
 e^{MR}\frac{(r/R)^{N+1}}{1-r/R}.
\]
For rational $G=P/Q$, the positive lower bound
$|Q_0|-\sum_{j\ge1}|Q_j|R^j$ gives $M$ directly.
These estimates, with outward interval operations, enclose the
root transports, the intertwiner and the four-equation real
system. Existence of the exact intertwiner is supplied by
\cref{night:P1.15}; an invertible square subsystem encloses it.
A matrix inverse on a rectangle and differentiation of
$Mv=b$ give
\[
 v_s=M^{-1}(b_s-M_sv),\qquad
 v_{st}=M^{-1}(b_{st}-M_{st}v-M_sv_t-M_tv_s).
\]
Together with $Y_x=YG$, $Y_y=iYG$ and
$Y_{xx}=Y(G^2+G')$, these enclose the real Hessian of $h_b$
on whole rectangles. For the value on a rectangle of
half-diagonal $\rho$, an entrywise bound $|G|\le\mathcal A$
gives
$|Y(z)-Y(z_0)|\le|Y(z_0)|_+(e^{\rho\mathcal A}-I)$.
All inputs to these enclosures are the displayed rational
connection, its gauges, the integral monodromies and the specified
normalization.

At infinity the deck action sends the normalized coefficient
pair to $(\zeta^8v_2,\zeta^4v_1)$, with
$\zeta=e^{2\pi i/5}$. Its value at the root point is therefore
zero. The height there is
$h_b(\infty)=-c_0\operatorname{Re}(F_\infty)_{03}$,
and the same deck action makes its gradient zero.
In the scaled coordinate $w=a^{1/5}u$, so $\lambda=aw^{-5}$,
the enclosures on the closed square $[-1/10,1/10]^2$ give
\[
 \min(h_{xx},h_{yy})-|h_{xy}|>\frac14.
\]
The real Hessian is consequently at least $I_2/4$ there.
Integrating along the segment from zero gives
\[
 h_b(aw^{-5})-h_b(\infty)\ge\frac18|w|^2.
\]
The normalized endpoint enclosure has radius below $4\cdot10^{-76}$
and lies between the negatives of the two endpoints in the statement.
\hypertarget{tag-03KY}{}\label{layout:V:fresh-minimum}

For the rest of the base use $v=a^{-1/4}u_0$, so
$\lambda=av^4$, and put $R=49/50$, $S=103/100$.
The first-quadrant $w$-disc $|w|\le R$ covers
$|\lambda|\ge aR^{-5}$ under its finite deck maps; the
first-quadrant $v$-disc covers $|\lambda|\le aR^4$.
The collar $R\le|w|\le S$ covers the intervening annulus,
since $R^4\ge S^{-5}$.
The cusp coordinate is $q=\lambda/\lambda_*-1=-w^{-5}-1$.
The same enclosure formulas on rational rectangular partitions
give the following strict lower bounds away from the central
square and the inner cusp disc:
\[
\begin{array}{c|c}
\text{region}&\text{lower bound for }h_b\\\hline
\text{infinity-root disc}&-1.094216\\
\text{zero-root disc}&-1.083056\\
\text{cusp annulus }1/100\le|q|\le1/4&-1.083777\\
\text{collar}&-1.030556
\end{array}
\]
Every endpoint in this table is an exact terminating decimal
strictly greater than $h_b(\infty)$. Rectangles meeting a
boundary are kept closed. Those lying in another of the stated
regions are handled there; those outside a specified radius are
excluded by their exact squared-distance inequality. The deck
maps and conjugation give the remaining sectors. Thus the
bounds cover all points outside $|q|<1/100$ and the central
square.

For the inner cusp disc, the exact gauge is
$D_\lambda=\operatorname{diag}(1,\lambda^2,\lambda,1,\lambda^{-1},\lambda^{-2})$.
Its connection is $R_c/q+A_0/(1+q)$ with $R_c^2=0$.
Writing $YD_\lambda=H_c(I+R_c\log q)P(q)$, $P(0)=I$,
the coefficient recurrence is
\[
 (nI-\operatorname{ad})P_n
 =\sum_{j<n}(-1)^{n-1-j}P_jA_0,
 \qquad\operatorname{ad}(Q)=QR_c-R_cQ.
\]
Here $\operatorname{ad}^3=0$, so the inverse is
$n^{-1}(I+\operatorname{ad}/n+\operatorname{ad}^2/n^2)$.
A scalar majorant $(1-|q|)^{-\gamma}$ gives the same
Cauchy tail control. Continuation to the positive cusp anchor
fixes $H_c$ and the logarithm branch. In the unimodular marking
$J$ with $J^{-1}NJ=\left(\begin{smallmatrix}0&I\\0&0\end{smallmatrix}\right)$,
the lower graph block is invertible throughout the closed
inner disc. The marking used for this anchor and the initial
cusp marking are related by $N=AN'A^{-1}$; their change has
block form $\left(\begin{smallmatrix}D&E\\0&D\end{smallmatrix}\right)$
and sends $Z$ to $DZD^{-1}+ED^{-1}$. It preserves the
real-intersection criterion.
\hypertarget{tag-03L0}{}\label{layout:V:fresh-marking}

At the parameter $\xi=ic_*$ write the resulting graph as
$Z(q)=\log q\,I_3/(2\pi i)+C(q)$, and put
$y_0=\log100/(2\pi)$ and $t=-\log|q|/(2\pi)-y_0\ge0$.
The coefficients, in increasing degree, of
$\det(tI_3+y_0I_3+\operatorname{Im}C(q))$ are bounded below
on $|q|\le1/100$ by
$3/2,9/2,15/4,1$, respectively.
The determinant is therefore positive for every nonzero $q$
in that disc. A graph vector $(Zv,v)$ is real precisely when
$v$ is real and $\operatorname{Im}Z\,v=0$. Thus the
$c_*$ plane is transverse there. The height at $q=1/100$
lies between $-0.451$ and $-0.450$; connectedness and
\eqref{maps:eq:F.22} give
$h_b(q)>h_b(\infty)$ on the entire punctured inner disc.

The strict central Hessian inequality, the four outer bounds and
the inner cusp inequality prove that infinity is the unique
minimum. The exact domain now follows from
\cref{night:P1.8}.
\hypertarget{tag-03L1}{}\label{layout:V:fresh-splitting}
\end{proof}
